%% OmniLaTeX Citation Style Library — Example
%%
%% This example demonstrates how to use \citationstyle{} to switch between
%% pre-configured bibliography styles.  Compile with:
%%   latexmk -xelatex main.tex
%%
%% Available styles: ieee, acm, apa, chicago, nature, science, harvard,
%%                   vancouver, mla

\documentclass[
    language=english,
    doctype=article,
    institution=none,
]{../../omnilatex}

\RequirePackage{omnilatex-citations}

%% Choose ONE citation style.  Uncomment the line you need:
\citationstyle{ieee}
% \citationstyle{acm}
% \citationstyle{apa}
% \citationstyle{chicago}
% \citationstyle{nature}
% \citationstyle{science}
% \citationstyle{harvard}
% \citationstyle{vancouver}
% \citationstyle{mla}

\addbibresource{refs.bib}

\begin{document}

\maketitle

\tableofcontents

\section{Introduction}
Citation styles vary across academic disciplines.  OmniLaTeX provides
pre-configured wrappers for the most common publishers so you can switch
styles with a single command.

As demonstrated by Knuth~\cite{knuth1984}, literate programming bridges
the gap between human-readable documentation and machine-executable code.
The BibTeX system~\cite{lamport1994} and its successor
biblatex~\cite{kime2020} provide powerful reference management.

Section~\ref{sec:methods} describes the methodology, which builds on
foundational work in software engineering~\cite{brooks1995}.

\section{Methods}
\label{sec:methods}
Our approach follows the structured methodology outlined in the
classic text by Brooks~\cite{brooks1995}, adapted for modern
open-source toolchains~\cite{kime2020}.

\section{Conclusion}
Switching citation styles is straightforward: change the
\texttt{\textbackslash citationstyle\{...\}} argument in the preamble
and recompile.

\printbibliography

\end{document}
